%!TEX root = ../main.tex
% Abstract in English (~250 words). Keywords are set in main.tex
% with \keywords{...}. Must be a faithful translation of the Resumen.

\begin{abstract}

This project was born as a nationwide network, with room to expand, to
connect learning and knowledge, elevating the tradition of academic
support in Argentine homes and giving it new tools. The initiative seeks
to make it easy for any student, from primary school to university, to
find the right academic support, allowing knowledge to flow without
geographical or physical barriers.

Today, the private tutoring market is highly fragmented. Although
relying on personal networks builds trust, its reach is very limited and
confined to the immediate geographical surroundings. The project
centralizes the offer in a digital ecosystem where search, booking,
virtual lessons, and payments coexist in a single place. This
professionalizes teaching work and offers families a transparent
environment to improve school results.

A central pillar is the trust model that safeguards students, especially
minors. Each user's identity is verified by reading their national
identity document on the platform's own server, tutors' academic
credentials are reviewed manually, and anyone who wants to teach minors
must submit a valid criminal record certificate. During every lesson, a
content classifier running in the browser acts as an emergency cut-off:
if a minor is present, it ends the session immediately. Minors have
their own account, but payments, authorizations, and reports are handled
by their responsible adult.

Artificial intelligence is applied where it adds concrete value: a
semantic matching engine compares what the student is looking for with
the topics each tutor teaches, and an automatic lesson summary is
offered as an optional add-on, only between adults and with the consent
of both. By valuing tutors' work, formalizing the sector, and adding
dynamic reputation and integrated virtual classrooms, the project
positions itself as an infrastructure that connects students with
qualified educators, overcoming geographical barriers and providing the
latter with a source of income. The impact of this solution will be
observed in the efficiency of the matching engine, the stability of
remote connections, and user satisfaction.

\end{abstract}
